% p.24 — الباب الثالث العشر في اعان يعين (اجوف في افعل)
\mspage{24}
\bab{\com{الباب الثالث العشر}{كذا في الأصل؛ ولعل الصواب: الباب الثالث عشر} في اعان يعين \qquad أَجْوَ\unclear{فُ} فِي أَفْعَلَ}
% Layout: six tall columns (right-to-left), each a stack of boxed blocks of unequal height.
% Column 1 is the summary column: the باب name (أَفْعَلَ) in its own box, then the ten
% underlying (pre-iʿlāl) forms أَعْوَنَ … لَا تُعْوِنْ, then the ten actual forms أَعَانَ … لَا تُعِنْ.
% In columns 2–6 every sub-heading box holds the heading and, beneath it, the underlying form
% (e.g. «ماضي معروف / أَعْوَنَ»); the regular paradigm follows outside the box.
% Row alignment is approximate (all columns run to 21 lines here); \cline marks the printed box
% boundaries inside each column.
% The chapter title is printed in two parts on one line: the باب title at the right, and the
% vocalised gloss «أَجْوَفُ فِي أَفْعَلَ» at the left.
\begin{sarf}{|c|c|c|c|c|c|}
أَفْعَلَ & \thead{ماضي معروف} & تُعِيْنُ & أُعِيْنَ & يُعَنَّ & أَعِيْنُوْا\\\cline{1-1}
أَعْوَنَ & أَعْوَنَ & تُعِيْنَانِ & أُعِيْنَا & تُعَانُ & أَعِيْنِيْ\\\cline{2-2}
يُعْوِنُ & أَعَانَ & يُعِنَّ & أُعِيْنُوْا & تُعَانَانِ & أَعِيْنَا\\
إِعْوَانًا & أَعَانَا & تُعِيْنُ & أُعِيْنَتْ & تُعَانُوْنَ & أَعِنَّ\\\cline{6-6}
مُعْوِنٌ & أَعَانُوْا & تُعِيْنَانِ & أُعِيْنَتَا & تُعَانِيْنَ & \thead{نهي حاضر}\\
أُعْوِنَ & أَعَانَتْ & تُعِيْنُوْنَ & أُعِنَّ & تُعَانَانِ & لَا تُعْوِنْ\\\cline{6-6}
يُعْوَنُ & أَعَانَتَا & تُعِيْنِيْنَ & أُعِنْتَ & تُعَنَّ & لَا تُعِنْ\\
إِعْوَانًا & أَعَنَّ & تُعِيْنَانِ & أُعِنْتُمَا & أُعَانُ & لَا تُعِيْنَا\\
مُعْوَنٌ & أَعَنْتَ & تُعِنَّ & أُعِنْتُمْ & نُعَانُ & لَا تُعِيْنُوْا\\\cline{5-5}
أَعْوِنْ & أَعَنْتُمَا & أُعِيْنُ & أُعِنْتِ & \thead{اسم مفعول} & لَا تُعِيْنِيْ\\
لَا تُعْوِنْ & أَعَنْتُمْ & نُعِيْنُ & أُعِنْتُمَا & مُعْوَنٌ & لَا تُعِيْنَا\\\cline{1-1}\cline{3-3}\cline{5-5}
أَعَانَ & أَعَنْتِ & \thead{اسم فاعل} & أُعِنْتُنَّ & مُعَانٌ & لَا تُعِنَّ\\\cline{6-6}
يُعِيْنُ & أَعَنْتُمَا & مُعْوِنٌ & أُعِنْتُ & مُعَانَانِ & \thead{مضارع بلم}\\\cline{3-3}
إِعَانَةً & أَعَنْتُنَّ & مُعِيْنٌ & أُعِنَّا & مُعَانُوْنَ & لَمْ يُعْوِنْ\\\cline{4-4}\cline{6-6}
مُعِيْنٌ & أَعَنْتُ & مُعِيْنَانِ & \thead{مضارع مجهول} & مُعَانَةٌ & لَمْ يُعِنْ\\
أُعِيْنَ & أَعَنَّا & مُعِيْنُوْنَ & يُعْوَنُ & مُعَانَتَانِ & لَمْ يُعِيْنَا\\\cline{2-2}\cline{4-4}
يُعَانُ & \thead{مضارع معروف} & مُعِيْنَةٌ & يُعَانُ & مُعَانَاتٌ & لَمْ يُعِيْنُوْا\\\cline{5-5}
إِعَانَةً & يُعْوِنُ & مُعِيْنَتَانِ & يُعَانَانِ & \thead{أمر حاضر} & لَمْ تُعِنْ\\\cline{2-2}
مُعَانٌ & يُعِيْنُ & مُعِيْنَاتٌ & يُعَانُوْنَ & أَعْوِنْ & لَمْ تُعِيْنَا\\\cline{3-3}\cline{5-5}
أَعِنْ & يُعِيْنَانِ & \thead{ماضي مجهول} & تُعَانُ & أَعِنْ & لَمْ يُعِنَّ\\
لَا تُعِنْ & يُعِيْنُوْنَ & أُعْوِنَ & تُعَانَانِ & أَعِيْنَا & لَمْ تُعِنْ\\
\end{sarf}
